\section{Data and methodology}
\label{sec:method}

A carrier from an accelerating transmitter drifts in sky frequency, and over
a long track its power spreads across many channels. The analysis therefore
\emph{de-drifts} the data, shifting each integration's spectrum to follow a
trial drift rate before adding the integrations together. This section
describes that search, how its output is turned into a power, and the three
steps through which every event then passes.

\subsection{The de-drifted search statistic}
\label{sec:statistic}

The search statistic is the largest stacked signal-to-noise ratio anywhere on
the channel$\times$drift grid of a window's dynamic spectrum,
%
\begin{equation}
T(\mathbf{x})=\max_{\dot\nu,\;\nu}\;
\frac{\sum_{t} I\!\left(\mathbf{x},t,\nu+\dot\nu t\right)
       \big/\sigma^{2}(t,\nu+\dot\nu t)}
     {\left[\sum_{t}1\big/\sigma^{2}(t,\nu+\dot\nu t)\right]^{1/2}},
\label{eq:tstar}
\end{equation}
%
where $I(\mathbf{x},t,\nu)$ is the spectrum at sky position $\mathbf{x}$ in
integration $t$ and channel $\nu$, $\dot\nu$ is the trial drift rate, and
$\sigma(t,\nu)$ is the noise scale, both evaluated at the channel the track
occupies in that integration. $T_\star$ denotes the value at the star. The
time $t$ runs from the first integration, so the frequency that labels a
track is the one a carrier would occupy at the start of the block, and a
repeat epoch is read at the cell the same convention fixes
(\S\ref{sec:etacrv}).

Equation~\ref{eq:tstar} is the matched filter, and $T$ an absolute
signal-to-noise ratio, only if the noise is independent between integrations
and channels. The analysis does not assume that. The trigger is instead set
against the distribution $T$ actually takes in the data, measured in every
window at \NCtrl{} control positions on the same visibilities and through the
same stack. A value of $T$ is then converted into a power with injected
carriers recovered through the unchanged pipeline (\S\ref{sec:pninetysel}).
The controls lie, from a fixed seed, in an annulus set by each block's
synthesised and primary beams. The noise scale $\sigma$ is their median
absolute deviation at each integration and channel, scaled to a
Gaussian-equivalent standard deviation, and the statistic is evaluated in the
same way at the star and at each control. The field read over is the whole
primary beam, that is every \texttt{FIELD} entry in which the star's
primary-beam response reaches \BgFieldMinResp.

Spectra are extracted at the star's apparent position, never from images. For
each integration and native channel the pipeline forms the weighted vector
average of the phase-shifted visibilities in total intensity, $I=(XX+YY)/2$.
That is the dirty-map value at the star, an amplitude and not a deconvolved
flux density, so it is not a point-source measurement where the field holds
extended emission. The correction for the star's offset from the phase
centre is measured with injected carriers (Appendix~\ref{app:pbresp}).
The deposited extractor propagates each position to the observation epoch
for proper motion and annual parallax. Some products were made by the first
forms of the extractor, which propagated proper motion only, and for those
windows the parallax loss is corrected window by window in
\S\ref{sec:pninetysel}.
Orbital motion is bounded for every component of a multiple in
Appendix~\ref{app:pbresp}.
%
% FLAGS SENTENCE (one sentence; every number from reflagdecl_r29.py, round 2900)
The data are the archive's deliveries. Of the \RflNEbCoreAll{} blocks that
set a quoted limit or hold a crossing, \RflNEbCore{} were recalibrated with
ALMA's final flags. They are not the same blocks as the \ZcNRecEB{} that
the recurrence test reads. The other \RflNEbNot{} keep the delivered
calibration, \RflNNotLim{} of them setting a limit and \RflNNotCross{}
holding a crossing. The change that the recalibration makes at the control
positions bounds the chance expectation of this mixed calibration at
\RflCeLo--\RflCeHi.

Every delivered channel is searched, without regridding, deconvolution or
edge trimming. A per-integration median along frequency removes the
continuum, and none is taken along time, so no feature narrow against the
band can be removed. The delivered channels are smoothed, so a correction
computed from the correlator's own channel shape turns flux density times
channel width into a radiated power (Appendix~\ref{app:conventions}).

The search itself runs on the topocentric sky frequencies ALMA delivered, and
every drift rate quoted is the apparent one at the telescope. Two later steps
need the star's own frame: comparing a crossing with the line catalogue, and
reading a repeat epoch at the \emph{predicted cell}, the single channel and
drift rate where the discovery's own fit puts a carrier. Both correct each
block's frequency axis for the Earth's motion and then for the star's
systemic velocity. That is the SIMBAD radial velocity \citep{SIMBAD2000},
queried on \SimbadQueryDate, or zero where none exists, as it is,
\FrZeroVsysClause. Toward $\beta$~Pic the value of the ALMA CO analysis of
its disc \citep{Matra2017} is adopted.

Two power quantities follow from the statistic, and both are defined in
Table~\ref{tab:quantities}. The trigger level $P_{\rm trig}$ is an internal
pipeline level. The sensitivity $\mathrm{EIRP}_{90}$ is measured by injecting
synthetic carriers of known amplitude into the calibrated visibilities of
\SensNInj{} windows and recovering them through the unmodified pipeline.
Every limit quoted is the $\mathrm{EIRP}_{90}$ of its own window, given beside
its native channel width, except the labelled bare-threshold comparisons of
\S\ref{sec:intro}, and its error budget is Table~\ref{tab:p90budget}. The
amplitude scale is checked twice. An independent reduction of the same
AU~Mic visibilities is reproduced to \PxJanRatio\,$\pm$\,\PxJanRatioErr,
which tests the extraction and not ALMA's flux transfer. The quiescent flux
density of $\epsilon$~Eridani \citep{Burton2022} is recovered to
\PxEpsRatio\,$\pm$\,\PxEpsRatioErr, which measures the scale to \SbEpsPct{}
per cent. That is consistent with ALMA's own accuracy specification,
$\pm\BudComb$ per cent for the sample, which we adopt for the absolute scale
(Table~\ref{tab:p90budget}).

The drift range is bounded by an acceleration and not by a frequency, since
drift rate and line-of-sight acceleration are one quantity in two sets of
units,
%
\begin{equation}
\frac{\dot\nu}{\nu}=\frac{a_{\rm los}}{c}
\qquad\text{and so}\qquad
a_{\rm los}=\frac{c\,\dot\nu}{\nu},
\label{eq:accel}
\end{equation}
%
and a ceiling on $\dot\nu/\nu$ is then the same physical limit in every band.
The ceiling is \DriftCeilLo\,Hz\,s$^{-1}$\,GHz$^{-1}$, or
\AccelCeilMed\,m\,s$^{-2}$ by Equation~\ref{eq:accel}, in \AccelNDefault{} of
the \NWinA{} windows, and \DriftCeilHi\,Hz\,s$^{-1}$\,GHz$^{-1}$
(\AccelCeilMax\,m\,s$^{-2}$) on the \AccelNWide{} \AccelWideStar{} windows,
whose innermost planet demands more. The default is roughly the acceleration
of a transmitter orbiting a solar-type star at a few hundredths of an au. It
is the \citet{Mason2025} fiducial with a threefold margin, and
\AccelSheikhPct{} per cent of the \AccelSheikhNHz\,nHz that
\citet{Sheikh2019} recommend for a general search. The grid is uniform in
$\dot\nu$ and runs from the negative ceiling to the first trial at or beyond
the positive one (Appendix~\ref{sec:edge}), in steps of half a native channel
over the span of the integrations searched. A carrier midway between two
trials therefore loses \DriftGridSpanHalf{} of a channel over that span, and
at most \DriftGridFineHalf{} over the longest on-source time.

The survey does not constrain a transmitter whose apparent
acceleration exceeds that range, whether it is chirped faster (a chirp is
not distinguished from acceleration) or its orbital and
rotational accelerations together exceed it. Transiting TRAPPIST-1~b is the
most demanding known case. At the default ceiling \TrapPhaseHi{} per cent of
its orbital phase would lie out of reach, and the widened grid used on its
star matches its maximum acceleration only to within the $\pm\TrpAmaxPct$ per
cent to which its host's mass and orbit are known.

Searching over drift also only helps if the largest searched drift moves a
carrier across about one channel while the star is observed,
$\eta_{\rm drift}=|\dot\nu_{\max}|\tau_{\rm on}/\Delta\nu_{\rm ch}\gtrsim1$,
where $\tau_{\rm on}$ is the time on source. \EtaDriftAOk{} of the \NWinA{}
Class~A windows satisfy it. The other \EtaDriftALo{} fall below it because
their tracks are short, and the per-window catalogue carries
$\eta_{\rm drift}$ for every window.

\subsection{From a crossing to a confirmed signal}
\label{sec:summary}

Every window passes through three steps in the same order: the trigger, the
line attribution and the recurrence test. They yield in turn a
\emph{crossing}, an \emph{unattributed crossing} and, had one survived, a
\emph{confirmed signal}, each defined once in Table~\ref{tab:quantities}.
Every count below is read from the ledger of every crossing the search
produced, and is of primary crossings, one to a window.

\emph{The trigger.} The level set on $T_\star$ carries no tail probability.
$T$ has not been shown to follow any standard distribution under the null
hypothesis, and a window holds between \TuCellsLo{} and $\TuCellsHi$
channel$\times$drift cells, so one level buys very different protection in
different windows. Every cell that reaches the trigger is recorded, and not
only the largest, because a carrier can be hidden as effectively by a noise
excursion as by a line.

\emph{The background.} With no tail probability attached, a fixed threshold
is not a fixed false-alarm rate either. The number of crossings expected by chance is instead computed window
by window, from the rate at which that window's own control positions exceed
the trigger. This assumes that a star behaves like a control position in its
own field, which the reserved hold-out blocks test (\S\ref{sec:trigunif}).
The census's \CsNCrossQp{} crossings, all of them Class~A, are the baseline
every expectation is set against. Of these, \BkFlagNCrossOther{} lies in a
window whose control ensemble fails to behave like noise
(Appendix~\ref{sec:blockquality}), and it is counted but withheld from the
amplitude-dependent tests (Appendix~\ref{app:falsealarm}).

\emph{The line attribution.} At these frequencies many crossings will be
molecular lines, so each one is compared with a line catalogue. The
comparison is a single database query, made on \MkListDate{} and repeatable
as it stands. A crossing within $\pm\MaskHalfKms$\,km\,s$^{-1}$ of a
catalogued transition in the star's rest frame is \emph{attributed}. The
half-width follows the Keplerian line-of-sight velocities of gas surviving
beyond a few au, with a margin for the systemic velocity. One fixed list of
\FrNSpec{} species is searched, whatever a window was tuned to. It was
completed at that query, after the crossings had been examined
(Appendix~\ref{app:linemask}), and the chance comparison of
\S\ref{sec:trigunif} is therefore conditional on the mask as adopted. Of the
\CnCross{} crossings of the census and the archival tail, the query attributes
\FrAttrWFifty{} and leaves \FrUnattrWFifty. Halving the half-width to
$\pm25$\,km\,s$^{-1}$ gives \FrAttrWTwentyFive{} and \FrUnattrWTwentyFive{}
respectively.
An attribution is a coincidence with a catalogue, not an identification or a
disposal, so every crossing goes on to the recurrence test whichever side of
the mask it falls. The price of the mask is that the survey constrains
nothing inside it.

\emph{The recurrence test.} A repeat counts only if it comes from a
\emph{separated} block, which here means a block of another observing session
and more than \EpSep{} day from the discovery. The session matters as much as
the interval, because blocks of one session share calibrators, atmosphere,
correlator setup and interference. Two readings are made. The
\emph{line-recurrence check} reads the predicted cell, $\pm\RqDvLo$ to
$\pm\RqDvHi$\,km\,s$^{-1}$ wide, in a separated block, and the
\emph{trajectory comparison} asks whether one bounded trajectory inside the
drift grid joins two crossings. Table~\ref{tab:quantities} states what each
reading can establish, and Fig.~\ref{fig:waterfall}(b) shows the frequencies
at which a separated block exists at all.

Whether the chain finds what it is meant to find is tested with the injected
carriers of \S\ref{sec:pninetysel}, and with $\beta$~Pictoris CO, a real
feature at a known stellar position, which is recovered and localised in two
transitions and \CsBpCoBlocks{} execution blocks (\S\ref{sec:bpicval}).


\begin{table*}
\centering
\caption{Operational definitions, used in this sense throughout. A crossing is
a single-epoch event; a confirmed signal needs a second, separated epoch,
which no sensitivity supplies.}
\label{tab:quantities}
\small
\begin{tabular}{@{}p{0.145\textwidth}p{0.825\textwidth}@{}}
\hline
Term & Definition \\
\hline
Crossing &
A channel\,$\times$\,drift cell at the stellar position reaching
$T_\star\geq\StageNullTrig$ in one execution block; neighbouring channels
above the trigger belong to one crossing. The largest in a window is its
\emph{primary} crossing, any other a \emph{secondary} one. An event requiring
investigation; no claimed significance. \\
Unattributed crossing &
A crossing more than $\MaskHalfKms$\,km\,s$^{-1}$ in the star's frame from
every transition of the line catalogue. A target for follow-up; not a
detection. \\
Confirmed signal &
The same carrier found again in a separated block (\S\ref{sec:summary}),
either at the cell the discovery predicts in the star's frame (the
\emph{line-recurrence check}) or in an additional separated epoch at the
frequency an independently specified trajectory predicts, with adequate
sensitivity there. The check recovers only an emitter at rest in the star's
frame, as a line is, so a null there never excludes a drifting or re-tuned
carrier. A pair of crossings that one bounded trajectory inside the drift
grid joins is \emph{open}, not confirmed: the joining test can reject a pair
and cannot accept one (Table~\ref{tab:pairs}). For \ZcNoRepSys{} Class~A
systems this archive holds no separated block at the frequency, so this test
can confirm nothing there, whatever the sensitivity. \\
$P_{\rm trig}$ &
$4\pi d^{2}S\,\Delta\nu$, the power at distance $d$ of a carrier that makes
$T_\star$ reach $T_{\rm trig}=\StageNullTrig$ in one native channel of width
$\Delta\nu$: an internal pipeline level, not a sensitivity. \\
$\mathrm{EIRP}_{90}$ &
The power at which nine injected carriers in ten produce a crossing at the
star, before the spatial rank against the control positions is applied. This
is what the paper means by a limit. It is a measured single-epoch trigger
completeness and should not be read as a detection probability. One class factor, calibrated by injection in \SensNInj{}
windows, is applied to every window, interpolated within the injected
configurations for \NgNInterpWin{} windows and extrapolated beyond them for
\NgNExtrapWin{} (\S\ref{sec:pninetysel}); each injected window's own point
($\times\SensIndivFacLo$--$\times\SensIndivFacHi$) is in the data archive.
It excludes atmospheric decorrelation (Table~\ref{tab:p90budget}). \\
Class~A &
The primary experiment, defined by channel width: \NWinA{} fine-channel
windows, median \ChanAMedKHz\,kHz, toward \NSysA{} systems over
\UnionClassA\,GHz of union bandwidth, \EtaDriftAOk{} of which satisfy
$\eta_{\rm drift}\gtrsim1$ (\S\ref{sec:statistic}). \\
Class~B &
The secondary experiment: \NWinB{} coarse-channel windows, mode
\ChanBModeMHz\,MHz, constraining an unresolved narrowband spectral excess.
Never a carrier search: in \ScCoarseNoDrift{} of them no searched drift leaves
one channel, and in the other \EtaDriftBHi{} the whole searched range moves a
carrier \ScBhEtaLo--\ScBhEtaHi{} channels, which registers one but does not
discriminate it. \\
\hline
\end{tabular}
\end{table*}
